% ============================================================
%  ANALYSIS
% ============================================================
\section{Analysis}

\subsection{Why Does ReLU\texorpdfstring{$^2$}{2} Resist Replacement?}

The persistence of $\mathrm{ReLU}^2$ as the top-performing activation across all tested
alternatives warrants careful examination.
NanoChat's activation was not chosen arbitrarily: prior work \parencite{so2021primer} found that
$\mathrm{ReLU}^2$ induces a higher degree of activation sparsity than GeLU in autoregressive
language models, and this sparsity has been linked to more efficient gradient flow and improved
loss scaling.
The sinusoidal perturbations introduced by GeLUSine and GSP add frequency-domain oscillations
that may benefit vision tasks---where natural image statistics have strong periodicity---but
introduce unnecessary noise in the token-level statistical setting of language modelling.
The GMTU's Gaussian damping envelope effectively zeroes out large activations, which may suppress
the model's ability to represent high-magnitude, semantically salient tokens.
The Turbulent activation's batch-level normalisation step couples activations across the batch
dimension, creating implicit dependencies that conflict with NanoChat's causal masking.

The key lesson is that activation function design interacts tightly with the training objective
and data modality.
Evolutionary search optimising for cross-dataset performance on vision tasks does not reliably
produce activations suited to language modelling, even at small scale.

\subsection{MTP: Richer Representations at a Throughput Cost}

The MTP results reveal an interesting tension.
The improved ChatCORE-Categorical performance on ARC and MMLU is consistent with the hypothesis
of \textcite{gloeckle2024better} that multi-horizon prediction forces the model to build more
structured, fact-oriented representations: correctly predicting four tokens ahead requires the
model to model syntactic and semantic structure beyond the immediate next word.
This is reflected in the marginal improvement on knowledge-retrieval benchmarks.

However, the throughput penalty is substantial.
MTP's four output heads increase GPU memory consumption significantly, reducing the effective
batch size at fixed memory budget.
This is a well-understood trade-off: the original MTP paper validates the approach on models
trained with hundreds of billions of tokens, where the richer representations compound over a
much longer training horizon.
At 5{,}000 iterations on a 2--4 GPU node, the additional compute cost is not amortised
sufficiently, and the aggregate gain is marginal.
Whether MTP becomes beneficial at longer training schedules (e.g.\ 50{,}000 iterations) is an
interesting open question for future work.

\subsection{Attention Residuals: Scale Sensitivity}

The failure of AttnRes to improve NanoChat points to a scale-dependent property of the mechanism.
AttnRes is motivated by \emph{PreNorm dilution}: the progressive weakening of early-layer signals
as depth increases with uniform additive accumulation.
At depth 12, however, this phenomenon is substantially less severe than at the depths of 40--80
layers typical of frontier models.
The hidden-state growth $O(L)$ at $L=12$ is only 12$\times$, compared to $O(80)$ for a
70-billion parameter model; the signal-dilution problem that AttnRes was designed to solve
therefore barely manifests.

Meanwhile, the additional parameters of the per-layer pseudo-query vectors $w_l$ and the
attention computation over stored layer outputs introduce optimisation complexity that requires
more data and iterations to converge.
At 5{,}000 training steps, the AttnRes model appears to be in a less-converged state than the
standard residual baseline, which benefits from a well-understood, well-initialised optimisation
landscape.
These observations suggest that AttnRes is primarily beneficial for large, deep models where
PreNorm dilution is severe and where the additional parameters can be meaningfully learned from
vast data.

\subsection{Engram: Scale Mismatch and Task Conflict}

The Engram results illustrate a fundamental challenge in transferring memory-augmented mechanisms
across scales.
The original Engram paper \parencite{cheng2026conditional} demonstrates gains at 27 billion
parameters trained on 262 billion tokens---roughly 310$\times$ the parameter count and
52$\times$ the token budget of our experiments.
At this scale, n-gram hash collisions are rare relative to vocabulary size, the memory store is
large enough to cache a meaningful fraction of common linguistic patterns, and the gating
mechanism has sufficient training signal to learn when to suppress the memory contribution.

At depth-12 / 85M parameters / 5{,}000 steps, none of these conditions hold.
The pretraining BPB curves overlapping between Engram and baseline indicate that the memory
module is providing no useful signal---the model has not accumulated enough context to populate
the store meaningfully, and the gate has not learned to control its influence.

The post-SFT degradation on SpellingBee is particularly revealing.
SpellingBee requires fine-grained compositional generation over individual characters, a task
that depends on the model tracking precise token-level context.
The Engram module's n-gram priors encode coarse word-level associations; when the gate fails to
suppress this signal during character-level reasoning, it corrupts the fine-grained representation.
The SFT data mixture used is not large enough for the gate to learn this suppression behaviour.
A complementary explanation is that HumanEval---also degraded by Engram---requires similar
compositional precision for syntactically correct code generation.

The one positive signal is the modest improvement on categorical knowledge tasks (ARC, MMLU)
with Engram, consistent with the memory store surfacing common factual n-gram patterns from
pretraining.
This suggests a fundamental trade-off: Engram increases \emph{knowledge retrieval} at the cost
of \emph{compositional generation}, and at this scale the trade-off is unfavourable.

\subsection{FP8: Implicit Regularisation as an Efficiency Gain}

The consistent superiority of the FP8 baseline across all metrics is a notable secondary finding.
Across ChatCORE, SpellingBee, and HumanEval, the FP8 configuration improves over its
full-precision counterpart without any architectural change.
One plausible explanation is \textbf{quantisation noise as implicit regularisation}: the rounding
errors introduced by 8-bit arithmetic act as a stochastic perturbation of the gradient, similar
in spirit to dropout or weight noise.
In the data-limited regime of 5{,}000 steps, this noise may prevent the model from overfitting
to training-corpus artifacts.
An alternative explanation is that FP8's higher arithmetic throughput allows more tokens to be
processed per wall-clock second, providing an effective increase in training data.
Disentangling these two effects is an interesting direction for future investigation.

\subsection{Qualitative Output Examples}

To complement the quantitative results, we examine sample generations from the best-performing
model configuration (d12\_baseline\_fp8 after SFT) and compare them with the standard baseline
and the Engram model.
The FP8 baseline generates coherent, grammatically well-formed text and produces plausible
answers to factual questions.
The Engram model's outputs on SpellingBee-style prompts (e.g., ``How many r's are in the word
\emph{strawberry}?'') show hallucination of incorrect character counts, consistent with the
quantitative benchmark degradation.
MTP-trained model outputs on MMLU-style prompts show slightly more confident, correct
multiple-choice responses, consistent with its categorical score advantage.
